\subsubsection{Espace des trajectoires et mesure invariante entre feuillets}
\label{sec:medida-retorno-hojas}

La topologie toroïdale d'APP fixe le support des translations cardinales
et leurs classes d'enroulement. La correspondance entre feuillets procède d'
une autre opération sur ce même support : la sélection temporelle de l'
évaluation additive ou multiplicative. Son typage géométrique est
\[
 \eta_{\rm av}(\Sigma)=\mathscr H_{\rm ar},\qquad
 \eta_{\rm av}(\Pi)=\mathscr H_{\rm vol}.
\]
Le premier feuillet correspond à la lecture additive de frontière ; le second,
à la lecture multiplicative d'intérieur. Cette application transporte les étiquettes
de mode et conserve l'incidence \eqref{eq:fav}. Le carré du rapport
d'auto-échelle se détermine au moyen des trajectoires de cette incidence.

\Needspace{8\baselineskip}
\begin{definition}[Enveloppe symbolique à mémoire un]
Soit
\[
 X_{\rm av}=\{(s_k)_{k\in\mathbb Z}\in
 \{\mathscr H_{\rm ar},\mathscr H_{\rm vol}\}^{\mathbb Z}:
 (F_{\rm av})_{s_k,s_{k+1}}=1\text{ pour tout }k\}.
\]
Le décalage des indices agit sur cet espace. C'est le plus petit système
symbolique à mémoire un qui contient les suites modales et conserve
leurs couples consécutifs : il admet les deux transitions croisées et l'
auto-rétention volumétrique. La phase, l'orientation cardinale et les registres
de la trajectoire restent dans l'état TPK à partir duquel on obtient cette
enveloppe.
\end{definition}

\begin{proposition}[Pondération des trajectoires entre feuillets]
Dans l'ordre surfacique--volumétrique, la mesure invariante d'entropie maximale
de \(X_{\rm av}\) a pour matrice de transition et distribution stationnaire
\begin{equation}
 P_{\rm av}=\begin{pmatrix}0&1\\\varphi^{-2}&\varphi^{-1}\end{pmatrix},
 \qquad
 (\mu_{\rm ar},\mu_{\rm vol})
 =\frac{(1,\varphi^2)}{1+\varphi^2}.
 \label{eq:medida-hojas-exacta}
\end{equation}
En particulier, \(\mu_{\rm ar}/\mu_{\rm vol}=\varphi^{-2}\).
\end{proposition}
\begin{proof}
L'incidence déjà construite possède le vecteur propre positif \(r=(1,\varphi)^{\mathsf T}\) et la valeur propre \(\varphi\). La normalisation
\[
 (P_{\rm av})_{ij}
 =\frac{(F_{\rm av})_{ij}r_j}{\varphi r_i}
\]
donne la matrice indiquée. Ses lignes ont pour somme un par \(\varphi^{-2}+\varphi^{-1}=1\). La symétrie de \(F_{\rm av}\) implique que les poids proportionnels à \(r_i^2\) satisfont l'équilibre détaillé ; ainsi, \(\mu P_{\rm av}=\mu\). L'irréductibilité et l'auto-rétention volumique déterminent une unique distribution stationnaire.

Pour vérifier la maximalité, toute transition permise satisfait
\(\log(P_{\rm av})_{ij}=\log r_j-\log r_i-\log\varphi\).
Sous toute mesure invariante, les termes d'extrémité se compensent
en moyenne. Si \(\nu_i\) sont ses poids de sommet et
\(q_i(j)=\nu(s_1=j\mid s_0=i)\), pour \(\nu_i>0\), alors
\[
 h_\nu\leq H_\nu(s_1\mid s_0)
 =\log\varphi-\sum_{i:\nu_i>0}\nu_iD(q_i\Vert P_i)
 \leq\log\varphi,
\]
où \(D(q\Vert p)=\sum_jq_j\log(q_j/p_j)\) est l'entropie
relative sur les couples permis, et l'on adopte \(0\log0=0\).
La mesure markovienne de \eqref{eq:medida-hojas-exacta} atteint la borne.
L'égalité dans la première inégalité exige la propriété de Markov ;
l'égalité dans la seconde exige \(q_i=P_i\). La stationnarité et
l'irréductibilité fixent alors \(\mu\), ce qui établit l'
unicité. C'est la mesure de Parry de l'enveloppe définie.
\end{proof}

La fréquence \((1/3,2/3)\) compte les instants du calendrier primitif ;
\(\mu\) pondère les trajectoires de l'enveloppe. Son rapport quadratique
exprime le produit des amplitudes d'incidence d'entrée et de sortie.
La mémoire et les restrictions de phase du TPK conservent leur domaine
propre ; l'entropie de l'enveloppe correspond au système symbolique
spécifié dans cette définition.

\subsubsection{Opérateur de clôture et limite des retours marqués}
\label{sec:operador-retorno-marcado}

La coordonnée de clôture \(\pi\), générée par le lecteur régional, intervient après la construction de l'incidence modale. Soient \(e_{\rm ar}=(1,0)^{\mathsf T}\), \(e_{\rm vol}=(0,1)^{\mathsf T}\) et \(P_i=e_ie_i^{\mathsf T}\) les projecteurs de feuillet. Les conditions du lecteur sont la conservation des deux feuillets, une normalisation volumique unitaire et une marque surfacique égale à la coordonnée de clôture. Ces conditions déterminent l'opérateur positif
\begin{equation}
 D_\pi=\pi P_{\rm ar}+P_{\rm vol}
       =\begin{pmatrix}\pi&0\\0&1\end{pmatrix}.
 \label{eq:operador-marca-pi}
\end{equation}
La diagonalité exprime la conservation des feuillets et les deux entrées
expriment leurs normalisations ; l'unicité de
l'opérateur pour le lecteur fixé est ainsi explicite.

\Needspace{8\baselineskip}
\begin{proposition}[Rapport des retours de frontière]
Pour \(n\geq1\), l'évaluation marquée
\begin{equation}
 \Theta_n=
 \frac{\langle e_{\rm ar},D_\pi F_{\rm av}^{\,n}e_{\rm ar}\rangle}
 {\langle e_{\rm vol},F_{\rm av}^{\,n}e_{\rm vol}\rangle}
 =\pi\frac{F_{n-1}}{F_{n+1}}
 \label{eq:retorno-marcado-finito}
\end{equation}
satisfait
\begin{equation}
 \lim_{n\to\infty}\Theta_n
 =\pi\frac{\mu_{\rm ar}}{\mu_{\rm vol}}
 =\frac{\pi}{\varphi^2}.
 \label{eq:retorno-marcado-limite}
\end{equation}
\end{proposition}
\begin{proof}
Les entrées diagonales de \(F_{\rm av}^{\,n}\) comptent les chemins
de longueur \(n\) qui reviennent à leur feuillet initial. La formule des puissances
déjà démontrée fournit \(F_{n-1}\) et \(F_{n+1}\). L'opérateur
\(D_\pi\) multiplie par \(\pi\) le premier dénombrement. Enfin,
\(F_{n-1}/F_{n+1}\) est le produit de deux quotients consécutifs
qui convergent vers \(\varphi^{-1}\) ; sa limite est
\(\varphi^{-2}\). L'identité avec le rapport des poids découle de
\eqref{eq:medida-hojas-exacta}.
\end{proof}

L'opérateur de clôture n'apparaît qu'une seule fois, dans l'évaluation terminale du retour. Le transfert \(F_{\rm av}\) reste déterminé par le calendrier. Ainsi se compose l'information de deux lecteurs du même état : l'incidence modale détermine la pondération quadratique et la région de clôture apporte sa coordonnée à la frontière.

La composition conserve aussi le raffinement. Pour des cylindres positifs
\(I=[a,b]\) et \(J=[c,d]\), l'image du lecteur
\(\mathcal R(x,y)=x/y^2\) est l'intervalle
\[
 \mathcal Q(I,J)=\left[\frac{a}{d^2},\frac{b}{c^2}\right].
\]
La monotonie en chaque argument prouve que les cylindres emboîtés de
clôture et d'auto-échelle produisent des intervalles emboîtés du rapport ; la
positivité de leurs limites et la contraction de leurs diamètres déterminent
\(\pi/\varphi^2\). L'expression scalaire retient ainsi une chaîne
d'opérations compatible avec la profondeur, qui s'applique
ci-dessous dans ses deux orientations.
